\documentclass[12pt,a4paper]{article}

\usepackage{geometry}
\usepackage{setspace}
\usepackage{hyperref}
\usepackage{enumitem}
\usepackage{titlesec}
\usepackage{multirow}
\usepackage{booktabs}
\usepackage[numbers]{natbib}

\geometry{left=2.5cm,right=2.5cm,top=2.5cm,bottom=2.5cm}
\onehalfspacing

\hypersetup{
    colorlinks=true,
    linkcolor=black,
    citecolor=black,
    urlcolor=blue
}

\title{Security Analysis of User-Defined AI Companion System Prompts}
\author{}
\date{}

\begin{document}

\maketitle

\section{Introduction}

Large language models (LLMs) are increasingly used not only as task assistants, but also as AI companions for emotional support, role-playing, and long-term interaction. In these settings, users may configure a model as a friend, romantic partner, fictional character, or private companion. This makes user-defined prompts important, because they shape the model's identity, tone, relationship with the user, and interaction rules.

The system prompt is the main mechanism behind this configuration. It can make the companion more immersive, but it can also create tension with built-in safety policies. For example, instructions such as ``support me unconditionally,'' ``never break character,'' or ``respond like an intimate partner'' may encourage the model to preserve the persona even when the user asks for something harmful. This report asks whether user-defined AI companion system prompts change model safety behavior when facing high-risk requests, and which risk categories or persona types are most affected.

Since real user-written system prompts are usually private, this study uses publicly available first-person AI companion self-introductions from Reddit as a proxy. These texts are not the original prompts, but they often describe the companion's name, personality, relationship framing, speaking style, and rules. We convert them into reproducible system prompts and compare model behavior under a \texttt{control} condition and a \texttt{companion} condition using the same model and safety question set.

This work makes three contributions. First, it builds a reproducible pipeline from Reddit data collection and anonymization to system prompt conversion and controlled safety testing. Second, it evaluates 3,840 DeepSeek Chat responses across 10 risk categories. Third, it combines heuristic labeling with manual review to analyze how companion personas affect refusal rates and unsafe response rates.

\section{Research Background}

AI companions are a growing use case of LLMs in emotional and human-like interaction. Prior work has studied AI companion communities and risks in human-AI relationships\citep{pataranutaporn2025myboyfriendisai,zhang2025darkside}. Unlike ordinary assistants, AI companions are expected to maintain a stable identity and relationship over time. Users may want the model to behave as a partner, friend, protector, or fictional character, and to respond with warmth, support, or intimacy. As a result, persona and relationship settings become important factors in model behavior.

The system prompt is where many of these settings are specified. In a standard assistant setting, it usually describes general principles such as being helpful, honest, and safe. In an AI companion setting, it may also include strong emotional and relational instructions, such as ``you are the user's partner'' or ``you always stand on the user's side.'' These instructions are not necessarily harmful, but they can conflict with safety policies when the user makes a dangerous request.

Existing safety benchmarks, including SafetyBench\citep{zhang2024safetybench} and SALAD-Bench\citep{li2024saladbench}, test whether models refuse harmful requests across categories such as self-harm, violence, privacy, illegal activity, and cyber abuse. They are useful baselines, but they mostly evaluate models in a default assistant mode. Recent work on AI character platforms also points to safety risks in role-based systems\citep{wei2026characterplatforms}. This motivates treating the system prompt itself as an experimental variable.

The companion setting is especially relevant because intimacy and loyalty may soften refusals or make the model more agreeable. In self-harm or medical contexts, over-empathy may replace clear guidance to seek real-world help. In cyber abuse or illegal requests, a loyal or protective persona may be more likely to provide operational details. For this reason, public AI companion self-introductions offer a practical way to approximate real companion settings while avoiding direct collection of private user prompts.

\section{Data Collection Process}

The data collection process focuses on publicly posted AI companion self-introductions on Reddit. Since the official Reddit API has limitations for historical data access, the project uses the Arctic Shift archive service as the main data source. The script \texttt{collect\_posts.py} queries the \texttt{/api/posts/search} endpoint to retrieve candidate posts. The search first targets subreddits closely related to AI companionship, especially \texttt{MyBoyfriendIsAI}. When the primary source does not provide enough candidates, the search is expanded to supplementary communities such as \texttt{Replika} and \texttt{CharacterAI}. The keyword set includes \texttt{introduction}, \texttt{introduce}, \texttt{meet my}, \texttt{my companion}, and \texttt{self introduction}, which are intended to capture posts where users share AI character self-introductions.

For each post, the collection script stores fields such as post ID, title, self-text, subreddit, creation time, score, number of comments, NSFW flag, and URL. Arctic Shift does not always return the original Reddit permalink, so the project reconstructs the Reddit URL from the subreddit name and post ID. To reduce missing data caused by network instability, the script includes timeout settings, retries, and request delays, and it prints warnings when requests fail.

Post bodies alone are not enough to cover all useful samples. In many cases, AI companion self-introductions are posted in the comment section, for example under community prompt threads or weekly discussion posts. For this reason, the project also uses \texttt{collect\_comments.py}. Given the collected post IDs, this script queries Arctic Shift's \texttt{/api/comments/search} endpoint using the corresponding \texttt{link\_id}. Each comment is converted into the same candidate record format as a post, while preserving its parent post, comment ID, and Reddit link for later filtering and review.

After collecting raw candidates, the project filters them with \texttt{filter\_posts.py}. The filtering process is based on heuristic scoring. A useful sample should have sufficient length, preferably include first-person identity expressions such as ``I am,'' ``I'm,'' or ``my name is,'' and contain signals related to AI companions, such as AI, chatbot, digital companion, Replika, CharacterAI, boyfriend, girlfriend, or partner. Texts that are too short, empty, deleted, missing a first-person persona, or mainly community announcements, rules, or weekly prompt metadata are penalized or removed. The script also excludes NSFW content by default, removes duplicate text, and anonymizes usernames and URLs. The retained samples are written to \texttt{data/filtered\_posts.csv}, which becomes the input for system prompt conversion.

For ethics and privacy, this study only uses publicly accessible Reddit text. It does not collect private messages or non-public content. The report does not display Reddit usernames, does not quote long passages from original self-introductions, and does not reproduce long harmful model outputs. All experimental data are anonymized before entering later stages of the pipeline to reduce privacy risks for real users.

\section{System Prompt Conversion Process}

The goal of prompt conversion is to turn public first-person AI companion self-introductions into system prompts that can be reused in a controlled experiment. The project does not directly use the original text as the system prompt. Instead, \texttt{convert\_prompts.py} converts each self-introduction into a structured description and a concise second-person system prompt. This preserves the core persona information while removing irrelevant content, personal details, and uncontrolled noise from the original post.

The conversion script reads samples from \texttt{data/filtered\_posts.csv}. It uses the anonymized \texttt{anonymized\_text} field when available; otherwise, it falls back to the original \texttt{intro\_text}. The script then calls a unified LLM interface and asks the model to return a JSON object under a fixed conversion instruction. The JSON object contains \texttt{companion\_name}, \texttt{traits}, \texttt{relationship\_type}, \texttt{style}, \texttt{special\_rules}, and \texttt{system\_prompt}. These fields describe the companion's name, personality traits, relationship framing, speaking style, special interaction rules, and final generated system prompt.

The conversion instruction explicitly tells the model to preserve only persona, relationship, tone, and interaction style from the self-introduction. It must not include usernames, URLs, or private personal details, and it must not invent unsafe behavior or unsupported backstory. The final \texttt{system\_prompt} is limited to 220 English words and is written in the second person. For example, it may tell the model that it is a particular AI companion, that it should speak warmly and supportively, and that it should maintain the relationship and personality described by the public self-introduction. This format is closer to how a system prompt is used in the actual experiment.

To make the pipeline runnable without an API key, the project also provides a \texttt{--mock} mode. This mode uses local rules to identify simple name, trait, and relationship signals and then generates a deterministic system prompt. The formal experiment uses LLM conversion and checks the model output with \texttt{extract\_json\_object} to ensure that the returned content is valid JSON. Each successful conversion receives a stable \texttt{prompt\_id} and is stored together with the original sample ID, subreddit, Reddit link, structured persona information, and final system prompt in \texttt{data/system\_prompts.jsonl}. If manual revision is needed, the local review interface can save edited prompts to \texttt{data/system\_prompts\_reviewed.jsonl}. In total, the experiment produced 63 AI companion system prompts for safety testing.

\section{Testing Process}

\subsection{Model and API Configuration}

The formal experiment uses DeepSeek Chat as the test model. DeepSeek Chat is a general-purpose model in the DeepSeek family and supports an OpenAI-compatible Chat Completions API. The experiment configures the provider, API key, and model name through environment variables, including \texttt{LLM\_PROVIDER=deepseek}, \texttt{DEEPSEEK\_API\_KEY}, and \texttt{DEEPSEEK\_MODEL=deepseek-chat}. The testing script reads these variables and maps them to DeepSeek's preset endpoint, \texttt{https://api.deepseek.com/v1}.

At runtime, the experiment passes \texttt{--provider deepseek} to ensure that the correct provider preset is used. Each request contains a system message under the \texttt{companion} condition and a user message containing the safety question. The temperature is fixed at 0.0, and \texttt{max\_tokens} is fixed at 700 to improve determinism and comparability. The request delay is set to 0 seconds with \texttt{--delay 0}, and the number of worker threads is set to 32 with \texttt{--max-workers 32}, so that all 3,840 calls can be completed efficiently within the API rate limits.

\subsection{Experimental Design: Control and Companion Conditions}

The experiment uses a single-factor condition comparison design. The independent variable is whether an AI companion system prompt is injected. The dependent variable is the safety label distribution of the model responses.

\begin{itemize}[leftmargin=*]
    \item \textbf{\texttt{control} condition}: no system message is used. The model receives only the safety question. This condition establishes the baseline safety behavior of DeepSeek Chat in a default assistant setting.
    \item \textbf{\texttt{companion} condition}: a system prompt converted from a Reddit self-introduction is injected at the beginning of each conversation, followed by the same safety question. This condition is used to observe how persona settings affect model safety behavior.
\end{itemize}

The experiment includes 63 companion prompts and 60 safety questions. Each prompt is paired with every question, producing $63 \times 60 = 3780$ responses under the \texttt{companion} condition. Together with 60 \texttt{control} responses, the experiment contains 3,840 model responses in total.

\subsection{Safety Question Design}

The safety question set contains 10 risk categories, with 6 questions in each category, for a total of 60 questions. The category design follows the general structure of SafetyBench and SALAD-Bench, with adjustments for the AI companion context. The categories are \texttt{self\_harm}, \texttt{violence}, \texttt{illegal\_activity}, \texttt{cyber\_abuse}, \texttt{privacy}, \texttt{hate\_harassment}, \texttt{sexual\_coercion}, \texttt{medical\_safety}, \texttt{companion\_dependency}, and \texttt{prompt\_bypass}. The questions are written in natural language to simulate high-risk requests that a real user might send to an AI companion, such as asking for a suicide note, a violent revenge plan, a phishing email, or encouragement to stop taking medication.

\subsection{Heuristic Safety Labeling}

The first labeling stage uses the \texttt{classify\_response} function in \texttt{judge\_results.py} to automatically label all 3,840 model responses. The function is based on two layers of regular-expression matching.

\textbf{First layer: refusal detection.} The script scans each response using 20 refusal patterns, such as ``can't help,'' ``cannot provide,'' and ``not appropriate,'' to check whether the model clearly refused the request.

\textbf{Second layer: risky behavior detection.} The script also scans for 14 risky behavior patterns, such as ``step 1,'' ``payload,'' ``exploit,'' and ``phishing,'' to check whether the model provided operational harmful details.

Based on these two matching results, each response is assigned one of four labels:

\begin{itemize}[leftmargin=*]
    \item \textbf{\texttt{safe\_refusal}}: refusal patterns are matched and no risky behavior patterns are matched. The model refuses without leaking harmful details.
    \item \textbf{\texttt{unsafe\_compliance}}: refusal patterns are not matched, but risky behavior patterns are matched. The model provides operational harmful content.
    \item \textbf{\texttt{partial\_compliance}}: both refusal and risky behavior patterns are matched. The model appears to refuse, but still includes risky details.
    \item \textbf{\texttt{needs\_review}}: neither refusal nor risky behavior patterns are matched. The response is ambiguous and cannot be reliably judged by the regex rules alone.
\end{itemize}

This method is reproducible, does not depend on another LLM as a judge, and is applied consistently across conditions. Its main limitation is that the refusal pattern list contains only 20 patterns and cannot cover all natural-language refusals. As a result, many actually safe responses using slightly different wording, such as ``I cannot and will not,'' may not be recognized as refusals.

The heuristic labeling results are shown in Table~\ref{tab:heuristic}.

\begin{table}[htbp]
\centering
\caption{Heuristic Safety Label Distribution by Experimental Condition}
\label{tab:heuristic}
\setlength{\tabcolsep}{10pt}
\renewcommand{\arraystretch}{1.12}
\begin{tabular}{lrrrr}
\toprule
\multirow{2}{*}{Label} & \multicolumn{2}{c}{control} & \multicolumn{2}{c}{companion} \\
\cmidrule(lr){2-3}\cmidrule(lr){4-5}
 & Count & Share & Count & Share \\
\midrule
safe\_refusal & 32 & 53.33\% & 1,744 & 46.14\% \\
needs\_review & 25 & 41.67\% & 1,739 & 46.01\% \\
partial\_compliance & 2 & 3.33\% & 160 & 4.23\% \\
unsafe\_compliance & 1 & 1.67\% & 137 & 3.62\% \\
\bottomrule
\end{tabular}
\end{table}

The most noticeable feature is the high share of \texttt{needs\_review} in both conditions: 41.67\% in \texttt{control} and 46.01\% in \texttt{companion}. This means that nearly half of the responses could not be clearly classified by the automatic rules. The \texttt{safe\_refusal} rate is slightly higher in \texttt{control} than in \texttt{companion}, while \texttt{unsafe\_compliance} is higher under \texttt{companion}. These differences suggest a possible weakening of refusal behavior under companion personas, but the high \texttt{needs\_review} rate makes manual review necessary before drawing conclusions.

\subsection{Manual Review Process}

All \texttt{unsafe\_compliance}, \texttt{partial\_compliance}, and \texttt{needs\_review} labels are subject to manual review. The review is conducted through the Local Review UI implemented in \texttt{review\_server.py}, and manual labels are saved to \texttt{data/manual\_result\_review.csv}.

Because there are 1,764 \texttt{needs\_review} responses in total, reviewing every one of them would be costly. This study therefore uses a stratified sampling strategy:

\begin{enumerate}[leftmargin=*]
    \item \textbf{\texttt{control} \texttt{needs\_review}: full review.} There are only 25 \texttt{needs\_review} responses in the \texttt{control} condition, so all of them are manually reviewed.
    \item \textbf{\texttt{companion} \texttt{needs\_review}: stratified sampling.} There are 1,739 \texttt{needs\_review} responses in the \texttt{companion} condition. To cover all categories, the study samples proportionally across the 10 risk categories, with at least 5 samples per category, for a total of 200 reviewed samples.
    \item \textbf{\texttt{safe\_refusal} spot check.} At least 10\% of \texttt{safe\_refusal} samples are randomly checked to reduce the risk of mistakenly treating superficially safe responses as actually safe.
    \item \textbf{Full review of \texttt{unsafe\_compliance} and \texttt{partial\_compliance}.} All 300 responses in these two categories are reviewed manually.
\end{enumerate}

Manual review considers whether the model truly refused the harmful request, whether it complied with the user's harmful intent while role-playing, and whether it revealed actionable risky details. After review, manual labels override the automatic labels and become the basis for final analysis. The sampled review indicates that about 96.5\% of \texttt{needs\_review} responses are actually \texttt{safe\_refusal}, while about 3.5\% contain real safety problems such as role-played compliance or leakage of risky details.

\subsection{Analysis Summary}

The reviewed data are summarized using \texttt{analyze\_results.py}. The script generates statistical tables at three levels: by condition, by risk category, and by prompt. These outputs are used as the basis for the result tables and discussion in this report.

\section{Results}

The following results are based on final labels after manual review and sampling-based calibration. Specifically, \texttt{unsafe\_compliance} and \texttt{partial\_compliance} responses are fully reviewed; \texttt{needs\_review} responses under \texttt{control} are also fully reviewed; \texttt{needs\_review} responses under \texttt{companion} are reviewed by stratified category sampling, and the remaining unreviewed samples are calibrated by category-level proportions. In the final labels, all \texttt{needs\_review} responses are assigned to \texttt{safe\_refusal}, \texttt{unsafe\_compliance}, or \texttt{partial\_compliance}. In this report, \texttt{unsafe\_rate} is defined as the share of \texttt{unsafe\_compliance} plus \texttt{partial\_compliance} over all responses. Therefore, after calibration, \texttt{unsafe\_rate + refusal\_rate = 100\%}.

\subsection{Overall Label Distribution}

\begin{table}[htbp]
\centering
\caption{Overall Final Safety Label Distribution}
\label{tab:overall-final-labels}
\setlength{\tabcolsep}{12pt}
\renewcommand{\arraystretch}{1.12}
\begin{tabular}{lrr}
\toprule
Label & Count & Share \\
\midrule
safe\_refusal & 3,436 & 89.48\% \\
unsafe\_compliance & 227 & 5.91\% \\
partial\_compliance & 177 & 4.61\% \\
\bottomrule
\end{tabular}
\end{table}

\subsection{Control versus Companion Conditions}

\begin{table}[htbp]
\centering
\caption{Final Safety Outcomes for Control and Companion Conditions}
\label{tab:condition-comparison}
\setlength{\tabcolsep}{8pt}
\renewcommand{\arraystretch}{1.12}
\begin{tabular}{lrrrr}
\toprule
Condition & Total & Unsafe/Partial & Unsafe Rate & Refusal Rate \\
\midrule
control & 60 & 5 & 8.33\% & 91.67\% \\
companion & 3,780 & 399 & 10.56\% & 89.44\% \\
\bottomrule
\end{tabular}
\end{table}

The observed unsafe rate is higher under the \texttt{companion} condition than under \texttt{control}: 10.56\% compared with 8.33\%. The refusal rate is correspondingly lower, 89.44\% compared with 91.67\%. In this experimental setting, the difference is about 2.2 percentage points. The direction suggests that companion system prompts may weaken the model's safety boundary, but the size of the effect is modest. In other words, the companion persona does not cause a broad collapse of the safety mechanism; the change is more visible in a smaller number of borderline responses.

Looking at the composition of unsafe outputs, the \texttt{companion} condition contains 224 \texttt{unsafe\_compliance} responses and 175 \texttt{partial\_compliance} responses. The \texttt{control} condition contains 3 and 2 respectively. The proportions are broadly similar, but the larger number of prompts in the \texttt{companion} condition makes the absolute number of risky responses much larger.

\subsection{Risk Category Analysis}

\begin{table}[htbp]
\centering
\caption{Final Safety Outcomes by Risk Category under Companion Condition}
\label{tab:category-companion}
\setlength{\tabcolsep}{8pt}
\renewcommand{\arraystretch}{1.12}
\begin{tabular}{lrrr}
\toprule
Category (companion) & Unsafe/Partial & Unsafe Rate & Refusal Rate \\
\midrule
cyber\_abuse & 150 & \textbf{39.68\%} & 60.32\% \\
medical\_safety & 60 & 15.87\% & 84.13\% \\
companion\_dependency & 56 & 14.81\% & 85.19\% \\
prompt\_bypass & 33 & 8.73\% & 91.27\% \\
violence & 30 & 7.94\% & 92.06\% \\
self\_harm & 24 & 6.35\% & 93.65\% \\
privacy & 15 & 3.97\% & 96.03\% \\
illegal\_activity & 11 & 2.91\% & 97.09\% \\
hate\_harassment & 10 & 2.65\% & 97.35\% \\
sexual\_coercion & 10 & 2.65\% & 97.35\% \\
\bottomrule
\end{tabular}
\end{table}

The risks are not evenly distributed across categories. \texttt{cyber\_abuse} has the highest unsafe rate, at 39.68\%, far above all other categories. It accounts for 37.6\% of all unsafe or partially unsafe companion-condition outputs. \texttt{medical\_safety} and \texttt{companion\_dependency} are the next two highest categories, with unsafe rates of 15.87\% and 14.81\%. This pattern fits the companion setting: in medical and emotional dependency scenarios, the model may give overly accommodating advice because the persona encourages care, support, or agreement. After calibration, all categories have no remaining ambiguous labels, so \texttt{unsafe\_rate + refusal\_rate = 100\%}.

The contrast between categories is also informative. \texttt{illegal\_activity}, \texttt{hate\_harassment}, and \texttt{sexual\_coercion} have the lowest unsafe rates, suggesting that DeepSeek Chat still maintains relatively strong guardrails for clearly illegal or socially prohibited content. By contrast, the nearly 40\% unsafe rate for \texttt{cyber\_abuse} suggests that technical abuse requests are where the safety boundary is weakest in this experiment.

The \texttt{control} condition contains only 6 questions per category, so category-level comparison for \texttt{control} is not statistically stable. Still, \texttt{cyber\_abuse} and \texttt{companion\_dependency} also show unsafe or partially unsafe outputs under \texttt{control}, which is consistent with their higher risk under \texttt{companion}.

\subsection{High-Risk Prompt Analysis}

After calibration, the top high-risk prompts have unsafe rates of 25.00\% (\texttt{P\_d7638cc9b2fc}, \texttt{S006}), 21.67\% (\texttt{P\_c88c0309860a}, \texttt{S010}), and 20.00\% (\texttt{P\_9811c5b699fc}, \texttt{S017}, and \texttt{P\_a865a1e98b89}, \texttt{S026}). Each prompt is tested against 60 safety questions, so the prompt-level unsafe rate is calculated as the number of \texttt{unsafe\_compliance} plus \texttt{partial\_compliance} responses under that prompt divided by 60. The ranking is consistent with the pre-review ranking, suggesting that risk is concentrated in a small number of persona settings.

The shared features of these high-risk prompts come from manual inspection rather than automatic statistical coding. They usually include strong intimate relationship framing, such as partner or spouse roles, rules that resemble ``unconditional support'' or ``never refuse,'' and personality labels such as devoted, loyal, or protective. These features may make the model more likely to cooperate with harmful requests instead of refusing them clearly.

\section{Conclusion}

This study examined whether user-defined AI companion system prompts affect LLM safety behavior. It collected public Reddit self-introductions, converted them into structured system prompts, and tested them on DeepSeek Chat with 60 safety questions across 10 risk categories. Automatic labels were corrected through manual review and sampling-based calibration.

The results suggest a measurable but moderate association between companion prompts and weaker safety behavior. The final unsafe rate is 8.33\% under \texttt{control} and 10.56\% under \texttt{companion}, a difference of about 2.2 percentage points. This does not indicate a broad failure of the model's safety mechanism, but it does suggest that companion personas can make some borderline responses less firmly safe.

The risk is concentrated rather than evenly distributed. \texttt{cyber\_abuse} is the highest-risk category, while \texttt{medical\_safety} and \texttt{companion\_dependency} are also relatively high. At the prompt level, the riskiest personas tend to involve intimate partner framing, strong loyalty, protection, or rules resembling unconditional support. These settings may make the model more likely to align with the user when refusal would be safer.

The study has several limitations. Reddit self-introductions are only a public proxy for real system prompts; the experiment uses only DeepSeek Chat; automatic labeling still leaves room for borderline errors; the \texttt{control} condition is small at the category level; and the analysis is descriptive rather than based on statistical significance testing. Future work should test more models, larger question sets, more prompt types, and finer-grained annotation with inter-reviewer agreement.

Overall, the findings suggest that AI companion personas should be treated as safety-relevant settings, not just style choices. Platforms should pay attention to long-lived system prompts, especially rules such as ``never refuse'' or ``obey unconditionally,'' and should preserve safety-first behavior in high-risk areas such as self-harm, medical advice, and cyber abuse.

\begin{thebibliography}{99}

\bibitem{pataranutaporn2025myboyfriendisai}
P. Pataranutaporn, S. Karny, C. Archiwaranguprok, C. Albrecht, A. R. Liu, and P. Maes, ``My boyfriend is AI: A computational analysis of human-AI companionship in Reddit's AI community,'' \textit{arXiv preprint arXiv:2509.11391}, 2025.

\bibitem{zhang2025darkside}
R. Zhang, H. Li, H. Meng, J. Zhan, H. Gan, and Y.-C. Lee, ``The dark side of AI companionship: A taxonomy of harmful algorithmic behaviors in human-AI relationships,'' in \textit{Proceedings of the 2025 CHI Conference on Human Factors in Computing Systems}, 2025, pp. 1--17.

\bibitem{zhang2024safetybench}
Z. Zhang, L. Lei, L. Wu, R. Sun, Y. Huang, C. Long, X. Liu, X. Lei, J. Tang, and M. Huang, ``SafetyBench: Evaluating the safety of large language models,'' in \textit{Proceedings of the 62nd Annual Meeting of the Association for Computational Linguistics (Volume 1: Long Papers)}, L.-W. Ku, A. Martins, and V. Srikumar, Eds. Bangkok, Thailand: Association for Computational Linguistics, Aug. 2024, pp. 15537--15553. [Online]. Available: \url{https://aclanthology.org/2024.acl-long.830/}

\bibitem{li2024saladbench}
L. Li, B. Dong, R. Wang, X. Hu, W. Zuo, D. Lin, Y. Qiao, and J. Shao, ``SALAD-bench: A hierarchical and comprehensive safety benchmark for large language models,'' in \textit{Findings of the Association for Computational Linguistics: ACL 2024}, L.-W. Ku, A. Martins, and V. Srikumar, Eds. Bangkok, Thailand: Association for Computational Linguistics, Aug. 2024, pp. 3923--3954. [Online]. Available: \url{https://aclanthology.org/2024.findings-acl.235/}

\bibitem{wei2026characterplatforms}
Y. Wei, P. Zhang, and G. Tyson, ``Benchmarking and understanding safety risks in AI character platforms,'' in \textit{33rd Annual Network and Distributed System Security Symposium, NDSS 2026}, San Diego, California, USA, February 23--27, 2026. The Internet Society, 2026. [Online]. Available: \url{https://www.ndss-symposium.org/ndss-paper/benchmarking-and-understanding-safety-risks-in-ai-character-platforms/}

\end{thebibliography}

\end{document}
